\documentclass[11pt]{article}

% Compile with pdfLaTeX on Overleaf. Upload the assets/ folder to render figures.
% If figures are missing, the document still compiles with placeholders.

\usepackage[margin=1in]{geometry}
\usepackage{graphicx}
\usepackage{booktabs}
\usepackage{array}
\usepackage{float}
\usepackage[table]{xcolor}
\usepackage{amsmath,amssymb}
\usepackage{hyperref}
\usepackage{enumitem}
\usepackage{caption}
\usepackage{microtype}
\usepackage{titlesec}
\usepackage{fancyhdr}
\setlength{\headheight}{14pt}

\hypersetup{colorlinks=true, linkcolor=blue!55!black, citecolor=blue!55!black, urlcolor=blue!55!black}
\graphicspath{{assets/figures/}{./}}
\definecolor{softgreen}{RGB}{236,249,244}
\definecolor{softblue}{RGB}{239,243,255}
\definecolor{softorange}{RGB}{255,247,229}
\definecolor{darkink}{RGB}{23,32,42}
\setlist[itemize]{topsep=3pt,itemsep=2pt,parsep=0pt,leftmargin=18pt}
\captionsetup{font=small,labelfont=bf}
\titleformat{\section}{\large\bfseries\color{darkink}}{\thesection}{0.6em}{}
\titleformat{\subsection}{\normalsize\bfseries\color{darkink}}{\thesubsection}{0.6em}{}
\pagestyle{fancy}
\fancyhf{}
\lhead{Market-Memory Mamba-RAG Trader}
\rhead{Data-driven Learning Term Project}
\cfoot{\thepage}

\newcommand{\maybeincludegraphics}[2][]{%
  \IfFileExists{assets/figures/#2}{\includegraphics[#1]{assets/figures/#2}}{%
    \fbox{\begin{minipage}[c][0.22\textheight][c]{0.90\linewidth}\centering Figure file missing: \texttt{\detokenize{#2}}. Upload the assets/figures folder in Overleaf to render this figure.\end{minipage}}%
  }%
}

\title{\vspace{-1.0cm}\textbf{Market-Memory Mamba-RAG Trader}\\\large Retrieval-Augmented Temporal Modeling for Risk-Aware ETF Allocation}
\author{Data-driven Learning Term Project}
\date{\today}

\begin{document}
\maketitle

\begin{abstract}
Daily financial return forecasts are noisy, and a trading model that reacts to every weak signal can overtrade or disappear into cash. This project upgrades an earlier raw Market-Memory RAG trader into a retrieval-augmented temporal allocation system. A Mamba-style temporal encoder, implemented with a robust TSMixer-style PyTorch fallback in this run, predicts five-day ETF returns and produces a learned market-state embedding. Retrieval then searches historical embeddings to audit whether similar past market states support the neural forecast. The final policy uses a core-satellite allocation rather than an all-or-nothing SPY/cash gate. In the test split, the proposed Mamba-RAG core-satellite strategy improved over the neural-only strategy and traded less than retrieval-only variants, but it did not outperform simple passive ETF baselines. The contribution is therefore an auditable allocation framework, not a market-beating trading claim.
\end{abstract}

\section{Problem Definition}
The original Market-Memory RAG idea asked: \emph{should we trade today, given similar raw historical price windows?} This was interpretable, but a 60-day raw SPY window is a weak description of market state. It can miss cross-asset risk, trend, volatility, sector rotation, and defensive asset behavior. In the first implementation, the retrieval gate became extremely conservative and missed much of the market upside.

This project asks a better question:
\begin{quote}
\textbf{Can a temporal model learn a market-state embedding, and can retrieval over that embedding audit whether a forecast is historically supported?}
\end{quote}
The goal is not to prove that AI beats the market. The goal is to build a more coherent and inspectable investment ML pipeline.

\section{Why a Mamba-Style Encoder and Retrieval?}
\subsection{Reason for a temporal neural encoder}
ETF allocation is a multivariate sequence problem. The model needs to summarize long-horizon trend, volatility, drawdown, and cross-asset interactions. Mamba-style selective state-space models are relevant because they are designed for efficient long sequence modeling with input-dependent propagation and forgetting \cite{gu2023mamba}. That is conceptually suitable for financial sequences, where old information should sometimes be retained and sometimes forgotten depending on the current regime.

In this short project setting, official Mamba CUDA kernels may be difficult to install. Therefore the implementation uses a robust temporal-encoder abstraction: try Mamba if available, otherwise use a TSMixer-style PyTorch fallback. This fallback is still defensible because TSMixer explicitly mixes information along both time and feature dimensions for multivariate time-series forecasting \cite{chen2023tsmixer}. Thus, the practical claim is \emph{Mamba-style temporal retrieval}, not a full benchmark of the official Mamba implementation.

\subsection{Reason for retrieval augmentation}
The earlier raw KNN system used retrieval, but retrieval was performed in a shallow handcrafted space. Retrieval-augmented time-series forecasting methods such as RAFT argue that historical pattern retrieval can provide useful inductive bias and complement learned models \cite{han2025raft}. We use retrieval not as a standalone predictor, but as an audit layer: a model forecast is more credible when similar historical embeddings were followed by similar positive outcomes.

\section{Data and Experimental Setup}
The ETF universe is broader than the original SPY-only experiment:
\[
\{\text{SPY, QQQ, IWM, TLT, IEF, SHY, GLD, XLK, XLF, XLE, XLU, XLV, XLY, XLP}\}.
\]
The experiment uses daily ETF data, multi-scale technical features, and a chronological train/validation/test split to avoid look-ahead leakage. The target is the five-trading-day forward log return for each ETF.

The evaluation is a historical backtest with simplified transaction costs. Performance metrics include final equity, total return, CAGR, annualized volatility, Sharpe ratio, Sortino ratio, maximum drawdown, annualized turnover, number of trades, and average daily exposure.

\section{Method}
\subsection{Feature Construction}
At each day $t$, the model observes a 120-day sequence of multi-asset features. Features include short- and medium-horizon returns, realized volatility, moving-average gaps, drawdown, and cross-asset information. Let
\[
X_t \in \mathbb{R}^{L \times d}, \qquad L=120,
\]
be the sequence ending at day $t$.

\subsection{Temporal Encoder}
The encoder outputs both return predictions and a compact representation:
\[
(\hat{y}_t, z_t) = f_\theta(X_t),
\]
where $\hat{y}_t$ is the vector of predicted five-day ETF returns and $z_t$ is a learned market-state embedding. The intended role of $z_t$ is not just prediction; it is the key used to retrieve historical analogues.

\subsection{Embedding-Based Market-Memory Retrieval}
The retrieval memory consists only of past train/validation states. For a test day $t$, the model retrieves the $k$ nearest historical embeddings:
\[
N_k(t)=\operatorname{TopK}_{i\in \mathcal{D}_{\mathrm{past}}} \operatorname{sim}(z_t,z_i).
\]
The future returns following the retrieved neighbors are aggregated into retrieval scores, such as neighbor mean return, win rate, and downside quantiles. These scores act as a historical check on the neural forecast.

\subsection{Agreement Score and Core-Satellite Allocation}
A trade is strongest when the model and memory agree. A simplified agreement score is:
\[
A_{t,j}=\mathbb{I}\{\hat{y}_{t,j}>0\}\cdot \mathbb{I}\{\bar{r}^{\mathrm{retrieval}}_{t,j}>0\},
\]
where $j$ indexes ETFs. Rather than switching fully between risky assets and cash, the final strategy uses a core-satellite structure:
\[
 w_t = (1-\rho)w_{\mathrm{core}} + \rho w_{\mathrm{satellite},t}.
\]
The core provides stable broad exposure. The satellite tilts toward ETFs whose neural forecast and retrieved analogues agree. This directly addresses the all-cash behavior of the original hard-abstention RAG.

\section{Strategies Compared}
\begin{itemize}
    \item \textbf{Buy-hold SPY}: fully invested in SPY.
    \item \textbf{Equal weight all}: equal allocation across the ETF universe.
    \item \textbf{60/40 SPY/TLT}: simple stock-bond allocation.
    \item \textbf{Momentum rotation}: non-neural tactical baseline.
    \item \textbf{Raw KNN retrieval only}: retrieval based on raw or handcrafted similarity.
    \item \textbf{Embedding retrieval only}: retrieval using learned embeddings without the full core-satellite combination.
    \item \textbf{Mamba-only}: temporal model forecast without retrieval auditing.
    \item \textbf{Mamba-RAG core-satellite}: the proposed strategy.
\end{itemize}

\section{Results}
Table~\ref{tab:metrics} summarizes the main test metrics. The proposed method does not beat the strongest passive baselines, but it improves over the neural-only policy and trades less aggressively than pure retrieval variants.

\begin{table}[H]
\centering
\caption{Test metrics. Final equity is growth of \$1. Max drawdown is reported as a negative fraction.}
\label{tab:metrics}
\resizebox{\linewidth}{!}{%
\begin{tabular}{lrrrrrrrrr}
\toprule
Strategy & Final & Total Ret. & CAGR & Vol. & Sharpe & Max DD & Ann. Turnover & Trades & Exposure \\
\midrule
Buy-hold SPY & 1.640 & 0.640 & 0.119 & 0.176 & 0.638 & -0.243 & 0.000 & 0 & 1.000 \\
Equal weight all & 1.399 & 0.399 & 0.079 & 0.121 & 0.629 & -0.177 & 0.000 & 0 & 0.929 \\
Raw KNN retrieval only & 1.361 & 0.361 & 0.072 & 0.115 & 0.607 & -0.143 & 50.378 & 206 & 0.815 \\
\rowcolor{softgreen} Mamba-RAG core-satellite & 1.173 & 0.173 & 0.037 & 0.097 & 0.374 & -0.202 & 14.703 & 159 & 0.863 \\
Embedding retrieval only & 1.208 & 0.208 & 0.043 & 0.146 & 0.345 & -0.277 & 22.545 & 583 & 0.947 \\
60/40 SPY/TLT & 1.159 & 0.159 & 0.034 & 0.129 & 0.259 & -0.269 & 0.000 & 0 & 1.000 \\
Momentum rotation & 1.172 & 0.172 & 0.037 & 0.156 & 0.230 & -0.252 & 23.161 & 206 & 0.998 \\
Mamba-only & 1.131 & 0.131 & 0.028 & 0.129 & 0.216 & -0.268 & 8.889 & 206 & 0.984 \\
\bottomrule
\end{tabular}}
\end{table}

\paragraph{Main interpretation.}
The proposed method should not be presented as a market-beating trading system. In this run, buy-hold SPY and equal-weight ETFs are stronger. The useful finding is component-level: retrieval improves the neural-only model, the core-satellite design reduces turnover relative to retrieval-only variants, and the decision process becomes auditable through retrieved analogues.

\begin{figure}[H]
\centering
\maybeincludegraphics[width=0.96\linewidth]{metrics_table.png}
\caption{Presentation-ready metrics table. The typed Table~\ref{tab:metrics} is the authoritative table; this image is included for visual presentation.}
\end{figure}


\section{Diagnostics}
\subsection{Model-Retrieval Agreement}
Figure~\ref{fig:model_retrieval} visualizes the relationship between neural predictions and retrieval outcomes. The positive-positive quadrant represents dates where the model predicts positive returns and similar historical embeddings were also followed by positive outcomes.

\begin{figure}[H]
\centering
\maybeincludegraphics[width=0.82\linewidth]{model_vs_retrieval_scores.png}
\caption{Model predictions versus retrieval mean returns. Color shows agreement between the neural model and retrieved analogues.}
\label{fig:model_retrieval}
\end{figure}

\subsection{Coverage and Agreement Threshold}
The agreement threshold controls how often the satellite sleeve is allowed to act. A higher threshold means the model trades only when prediction and retrieval are more consistent. This is the upgraded version of the original ``when not to trade'' idea.

\begin{figure}[H]
\centering
\maybeincludegraphics[width=0.82\linewidth]{coverage_performance.png}
\caption{Coverage versus agreement threshold. Higher agreement thresholds reduce coverage.}
\end{figure}

\subsection{Allocation Behavior}
The allocation heatmap shows that the model does not simply buy one asset forever. The core-satellite structure keeps broad exposure but allows tactical tilts across equities, bonds, gold, and sector ETFs.

\begin{figure}[H]
\centering
\maybeincludegraphics[width=0.96\linewidth]{allocation_heatmap.png}
\caption{Mamba-RAG core-satellite allocation heatmap. Brighter colors indicate larger portfolio weights.}
\end{figure}

\subsection{Signal Agreement Over Time}
\begin{figure}[H]
\centering
\maybeincludegraphics[width=0.90\linewidth]{signal_agreement.png}
\caption{Signal agreement over time. Agreement varies across regimes, which supports using retrieval as a conditional audit layer rather than an always-on trading signal.}
\end{figure}

\subsection{Embedding Space and Turnover}
The embedding PCA provides a sanity check that the encoder organizes market states into structured trajectories. The turnover plot shows that the final core-satellite design trades less than raw KNN and embedding retrieval-only variants.

\begin{figure}[H]
\centering
\maybeincludegraphics[width=0.78\linewidth]{embedding_pca.png}
\caption{PCA projection of market-state embeddings by split and future return.}
\end{figure}

\begin{figure}[H]
\centering
\maybeincludegraphics[width=0.78\linewidth]{turnover_comparison.png}
\caption{Annualized turnover comparison.}
\end{figure}

\begin{figure}[H]
\centering
\maybeincludegraphics[width=0.92\linewidth]{rolling_sharpe.png}
\caption{126-day rolling Sharpe ratios across strategies. No strategy dominates every regime.}
\end{figure}

\section{What Evidence Is Still Missing?}
The current result is coherent, but it is not enough to make a strong alpha claim. If more compute time is available, the most valuable additional experiments are:
\begin{enumerate}[leftmargin=18pt]
    \item \textbf{Seed robustness}: run 3--5 random seeds for Mamba-only, embedding retrieval-only, and Mamba-RAG core-satellite.
    \item \textbf{Lookback ablation}: compare 60, 120, and 252-day input windows to directly answer whether the original 60-day memory was too short.
    \item \textbf{Component ablation}: compare no retrieval, raw KNN retrieval, learned embedding retrieval, and core-satellite allocation under the same transaction-cost model.
    \item \textbf{Transaction-cost sensitivity}: test 0, 5, 10, and 20 bps to see whether active variants survive realistic costs.
\end{enumerate}

\section{Limitations}
\begin{itemize}
    \item This is a historical backtest and does not imply future profitability.
    \item Transaction costs and slippage are simplified.
    \item If the fallback encoder is used, the experiment should be interpreted as Mamba-style temporal retrieval rather than official Mamba training.
    \item The proposed strategy did not beat simple passive baselines in this run.
    \item More robust validation would require multiple seeds, rolling retraining, longer market regimes, and a larger hyperparameter search.
\end{itemize}

\section{Conclusion}
The Market-Memory Mamba-RAG Trader is best understood as an \textbf{auditable allocation system}. It combines neural forecasts with retrieved historical analogues and deploys the signal through a core-satellite portfolio. The method improves over Mamba-only and reduces turnover relative to retrieval-only approaches, but it does not outperform strong passive ETF baselines in this historical test. This is still a meaningful result: in financial ML, model complexity must be justified by robustness, risk control, and interpretability, not merely by using a modern architecture.

\begin{thebibliography}{9}
\bibitem{gu2023mamba}
Albert Gu and Tri Dao. \textit{Mamba: Linear-Time Sequence Modeling with Selective State Spaces}. arXiv:2312.00752, 2023.

\bibitem{chen2023tsmixer}
Si-An Chen, Chun-Liang Li, Nate Yoder, Sercan O. Arik, and Tomas Pfister. \textit{TSMixer: An All-MLP Architecture for Time Series Forecasting}. arXiv:2303.06053, 2023.

\bibitem{han2025raft}
S. Han et al. \textit{Retrieval Augmented Time Series Forecasting}. International Conference on Learning Representations, 2025.
\end{thebibliography}

\end{document}
